% ===============================================================
% LH-05b: Lehrerhinweise - Reguläre Verben auf -er/-ir
% Abgeleitet aus LE-05b (Verlauf, Differenzierung, Fehler)
% ===============================================================

\documentclass[11pt, a4paper]{article}

\usepackage[utf8]{inputenc}
\usepackage[T1]{fontenc}
\usepackage[ngerman]{babel}
\usepackage{helvet}
\renewcommand{\familydefault}{\sfdefault}

\usepackage[margin=2cm]{geometry}
\usepackage{enumitem}
\usepackage{tabularx}
\usepackage{booktabs}
\usepackage{longtable}

\usepackage{xcolor}
\usepackage{tcolorbox}

\usepackage{fancyhdr}
\usepackage{lastpage}
\usepackage{amssymb}

% Farben
\definecolor{spanisch}{HTML}{c41e3a}
\definecolor{grau}{HTML}{888888}
\definecolor{hellgrau}{HTML}{f5f5f5}
\definecolor{basis}{HTML}{2a7a4b}
\definecolor{standard}{HTML}{2c5aa0}
\definecolor{erweiterung}{HTML}{7b1fa2}
\definecolor{loesung}{HTML}{c62828}

\newcommand{\fachfarbe}{spanisch}

% Variablen
\newcommand{\thema}{Reguläre Verben auf -er/-ir}
\newcommand{\sequenz}{05}
\newcommand{\block}{b}
\newcommand{\fach}{Spanisch}
\newcommand{\klasse}{7}
\newcommand{\dauer}{90 Min. (Doppelstunde)}

% Kopf-/Fußzeile
\pagestyle{fancy}
\fancyhf{}
\renewcommand{\headrulewidth}{0.4pt}
\renewcommand{\footrulewidth}{0.4pt}
\lhead{\fach{} -- Klasse \klasse}
\chead{\textbf{LH-\sequenz\block{} (Lehrerhinweise)}}
\rhead{}
\lfoot{}
\rfoot{Seite \thepage\ von \pageref{LastPage}}

% Differenzierungsmarker
\newcommand{\B}{\textcolor{basis}{\textbf{[B]}}}
\newcommand{\Std}{\textcolor{standard}{\textbf{[S]}}}
\newcommand{\E}{\textcolor{erweiterung}{\textbf{[E]}}}
\newcommand{\lsg}[1]{\textcolor{loesung}{\textbf{#1}}}
\newcommand{\es}[1]{\textcolor{\fachfarbe}{\textbf{#1}}}

% Boxen
\newtcolorbox{kurzplan}{
    colback=hellgrau,
    colframe=\fachfarbe,
    boxrule=1pt,
    arc=0pt,
    fonttitle=\bfseries,
    title=Kurzplan
}

\newtcolorbox{differenzierung}{
    colback=white,
    colframe=grau,
    boxrule=0.5pt,
    arc=0pt,
    fonttitle=\bfseries,
    title=Differenzierung
}

\newtcolorbox{fehlerbox}{
    colback=white,
    colframe=loesung,
    boxrule=0.5pt,
    arc=0pt,
    fonttitle=\bfseries,
    title=Häufige Fehler
}

\newtcolorbox{materialbox}{
    colback=white,
    colframe=\fachfarbe,
    boxrule=0.5pt,
    arc=0pt,
    fonttitle=\bfseries,
    title=Material-Checkliste
}

% ===============================================================
\begin{document}

\begin{center}
    {\Large\bfseries\textcolor{\fachfarbe}{Lehrerhinweise: \thema}}\\[0.3em]
    {\normalsize LH-\sequenz\block{} | \dauer{} | Std. 3-4 von 8}
\end{center}

\vspace{10pt}

% ===============================================================
% 1. KURZPLAN
% ===============================================================

\section*{1. Stundenverlauf (Kurzplan)}

\textbf{1. Stunde (45 Min.) -- Fokus: Verben auf -er}

\begin{tabularx}{\textwidth}{|c|l|X|l|}
\hline
\textbf{Zeit} & \textbf{Phase} & \textbf{Inhalt} & \textbf{Material} \\
\hline
0--5 & Einstieg & Bildimpuls: Freizeitaktivitäten & Bilder \\
\hline
5--10 & Aktivierung & Wiederholung -ar Endungen & Tafel \\
\hline
10--20 & Input & -er Verben einführen: comer, beber, leer, correr & Tafel, AB \\
\hline
20--30 & Übung 1 & Chorisches Sprechen: Konjugation -er & -- \\
\hline
30--40 & Übung 2 & EA: Lückentext AB-05b, Aufgabe 1 & AB-05b \\
\hline
40--45 & Sicherung & Lösungen besprechen, Merkkasten -er & AB-05b \\
\hline
\end{tabularx}

\vspace{1em}

\textbf{2. Stunde (45 Min.) -- Fokus: Verben auf -ir + Vergleich}

\begin{tabularx}{\textwidth}{|c|l|X|l|}
\hline
\textbf{Zeit} & \textbf{Phase} & \textbf{Inhalt} & \textbf{Material} \\
\hline
0--5 & Aktivierung & Quiz: -er Endungen abfragen & -- \\
\hline
5--15 & Input & -ir Verben: vivir, escribir, abrir. Vergleich mit -er & Tafel \\
\hline
15--25 & Übung 3 & PA: Sätze bilden (AB-05b, Aufgabe 3) & AB-05b \\
\hline
25--35 & Vertiefung & Vergleichstabelle -ar/-er/-ir (Aufgabe 4) & AB-05b \\
\hline
35--40 & Sicherung & Merkkasten -ir, Unterschiede zusammenfassen & AB-05b \\
\hline
40--45 & Abschluss & Zusammenfassung, HA & -- \\
\hline
\end{tabularx}

% ===============================================================
% 2. DIFFERENZIERUNG
% ===============================================================

\section*{2. Differenzierung}

\begin{differenzierung}
\textbf{Legende:} \B{} Basis (BR) | \Std{} Standard (MR) | \E{} Erweiterung (GY)

\vspace{0.5em}
\textbf{WICHTIG:} Diese Marker sind NUR für die Lehrkraft!
\end{differenzierung}

\vspace{1em}

\begin{tabularx}{\textwidth}{|l|X|X|X|}
\hline
& \B{} \textbf{Basis} & \Std{} \textbf{Standard} & \E{} \textbf{Erweiterung} \\
\hline
\textbf{Aufg. 1-2} & Endungstabelle sichtbar & Selbstständig & + Zusatzsätze \\
\hline
\textbf{Aufg. 3} & Wortschatz vorgegeben & Mit Hilfe & Freie Satzbildung \\
\hline
\textbf{Aufg. 4} & Tabelle vorausgefüllt & Vollständig & + Transfer auf neue Verben \\
\hline
\end{tabularx}

\vspace{1em}

\textbf{Konkrete Hinweise:}
\begin{itemize}
    \item \B{} LRS-SuS: Zeitzugabe (+10 Min.), Endungstabelle als Hilfe erlauben
    \item \B{} DaZ-SuS: Deutsche Bedeutungen zu jedem Verb notieren
    \item \Std{} Standardgruppe: Partnerarbeit bei Aufgabe 3
    \item \E{} Leistungsstarke: Zusätzliche Verben konjugieren (aprender, comprender, decidir)
\end{itemize}

% ===============================================================
% 3. HÄUFIGE FEHLER
% ===============================================================

\section*{3. Häufige Fehler}

\begin{fehlerbox}
\begin{tabularx}{\textwidth}{|l|X|X|}
\hline
\textbf{Fehler} & \textbf{Beispiel} & \textbf{Reaktion} \\
\hline
Verwechslung -emos/-imos & *``Nosotros comimos'' statt ``comemos'' & Farbcodierung: -er blau, -ir grün \\
\hline
Verwechslung -éis/-ís & *``Vosotros vivéis'' statt ``vivís'' & Akzentregel: -ís hat Akut auf i \\
\hline
Transfer von -ar & *``Yo como'' vs. *``Yo coma'' & Vergleichstabelle: -o bleibt gleich! \\
\hline
Aussprache ``leer'' & *[lia] statt [le.'er] & Vorsprechen: le-er (zwei Silben) \\
\hline
Stammfehler & *``viver'' statt ``vivir'' & Infinitiv an die Tafel, betonen \\
\hline
\end{tabularx}
\end{fehlerbox}

\vspace{1em}

\textbf{Merksatz für SuS:}

\begin{center}
\fbox{\parbox{0.8\textwidth}{\centering
\textbf{-er und -ir sind fast gleich!}\\
Nur bei \textit{nosotros} und \textit{vosotros} unterscheiden sie sich:\\
com\textcolor{standard}{\textbf{emos}}/com\textcolor{standard}{\textbf{éis}} vs. viv\textcolor{erweiterung}{\textbf{imos}}/viv\textcolor{erweiterung}{\textbf{ís}}
}}
\end{center}

% ===============================================================
% 4. MATERIAL-CHECKLISTE
% ===============================================================

\section*{4. Material-Checkliste}

\begin{materialbox}
\begin{itemize}[label=$\square$]
    \item AB-05b.pdf (24 Kopien)
    \item ML-05b.pdf (1x für Lehrkraft)
    \item Lehrwerk: Qué pasa 1, S. 44--46
    \item Bilder: Freizeitaktivitäten (essen, lesen, laufen, schreiben)
    \item Verbkarten für PA (Kopiervorlage, 12 Sets)
    \item Optional: LP-05b.html (Tablets/PCs)
\end{itemize}
\end{materialbox}

% ===============================================================
% 5. LÖSUNGEN
% ===============================================================

\section*{5. Lösungen AB-05b}

\textbf{Merkkasten 1 (-er):}
como, comes, come, comemos, coméis, comen

\textbf{Merkkasten 2 (-ir):}
vivo, vives, vive, vivimos, vivís, viven

\vspace{0.5em}

\textbf{Aufgabe 1:} (je 1P)
\begin{enumerate}[label=\alph*)]
    \item \lsg{como}
    \item \lsg{bebes}
    \item \lsg{lee}
    \item \lsg{corremos}
    \item \lsg{beben}
\end{enumerate}

\textbf{Aufgabe 2:} (je 1P)
\begin{enumerate}[label=\alph*)]
    \item \lsg{vivo}
    \item \lsg{escribes}
    \item \lsg{abre}
    \item \lsg{vivimos}
    \item \lsg{escribís}
\end{enumerate}

\textbf{Aufgabe 3:} (je 1P)
\begin{enumerate}[label=\alph*)]
    \item \lsg{Yo leo un libro en casa.}
    \item \lsg{Mi hermano corre en el parque.}
    \item \lsg{Nosotros comemos pizza con amigos.}
    \item \lsg{Tú vives en España.}
\end{enumerate}

\textbf{Aufgabe 4:} (4P Tabelle + 2P Erklärung)

Tabelle:
\begin{itemize}
    \item comer: como, comes, come, comemos, coméis, comen
    \item vivir: vivo, vives, vive, vivimos, vivís, viven
\end{itemize}

Erklärung: \lsg{Der Unterschied zwischen -er und -ir liegt nur bei nosotros (-emos vs. -imos) und vosotros (-éis vs. -ís).}

% ===============================================================
% 6. HAUSAUFGABE
% ===============================================================

\section*{6. Hausaufgabe}

\begin{itemize}
    \item Lehrwerk S. 45, Übung 4 (Verben konjugieren)
    \item Vokabeln lernen: comer, beber, leer, correr, vivir, escribir, abrir
    \item Optional: Lernpfad LP-05b.html bearbeiten
\end{itemize}

% ===============================================================
% 7. REFLEXIONSNOTIZEN
% ===============================================================

\section*{7. Reflexionsnotizen (nach Durchführung)}

\begin{tabularx}{\textwidth}{|l|X|}
\hline
\textbf{Was lief gut?} & \\[2em]
\hline
\textbf{Was lief nicht?} & \\[2em]
\hline
\textbf{Zeitmanagement} & \\[2em]
\hline
\textbf{Für nächstes Mal} & \\[2em]
\hline
\end{tabularx}

\end{document}
